% !TeX root = ../Project.tex
\section{Introduction}
\label{sec:introduction}

Proteins are fundamental components of living cells and drive nearly every biological process. Understanding how their sequence shapes structure and function could help us explain existing biology and design new proteins, with broad potential across science, medicine and engineering \citep{huang2016coming,kuhlman2019advances}. In only a few years, AlphaFold brought tremendous progress in protein structure prediction, while generative models moved beyond predicting natural proteins towards designing new ones. Models such as RFdiffusion generate protein structures around functional motifs, while ESM3 generates across sequence, structure and function \citep{jumper2021,watson2023rfdiffusion,hayes2025esm3}. Together, these developments have rapidly accelerated protein design.

All-atom modelling is particularly interesting because protein function depends not only on the overall backbone, but also on amino-acid identity, side-chain placement and the local atomic interactions involved in binding and catalysis. Protpardelle generates protein sequences and all-atom structures directly, while La-Proteina learns residue representations through reconstruction and uses its latents for all-atom generation \citep{chu2024allatom,la-proteina}. These objectives require representations to retain detailed atomic information. However, each residue is also shaped by its wider structural environment, including residues distant in sequence but close in space. Building on this progress, our aim is to learn representations that help us understand proteins and can eventually support the generation of new ones. The question is whether a representation can capture both local and broader structure in a way that transfers to biological tasks.

Joint-Embedding Predictive Architectures (JEPAs) offer an alternative to reconstruction-based learning by predicting representations in latent space rather than reconstructing the raw input. With the swift progress of JEPA-style learning in image and text domains, similar approaches have come within reach for applications in structural biology \citep{assran2023ijepa,balestriero2025lejepa}. In image-based I-JEPA, this produces representations of high-level structure rather than low-level pixel details \citep{assran2023ijepa}. I-JEPA shows that masking is central to this: target regions must contain meaningful information, while the visible context must remain informative enough to predict them. For proteins, different masking strategies could let us examine what residue information can be learned from the remaining structural context.

Current protein JEPAs operate on sequence or multimodal supervision rather than atomic structure \citep{ofer2026proteinjepa,ravideshik2026protjepa}. To our knowledge, this is the first study of JEPA latent prediction over residue-level all-atom structures, and a step towards an all-atom protein world model: rather than predicting 3D coordinates directly, the model predicts learned residue representations of hidden structure from visible context. While this thesis studies JEPA for representation learning, a fuller protein world model would also predict how protein states change through protein dynamics or generative state transitions \citep{maes2026leworldmodel}.

These directions lead to the central research question of this thesis:

\researchquestionbox{Can JEPA-style learning produce useful all-atom protein representations, and how do protein-specific design choices shape what they retain?}

\paragraph{Contributions.}
To answer this question, this thesis makes four contributions:
\begin{enumerate}[leftmargin=*,itemsep=4pt,topsep=4pt,parsep=0pt]
    \item \textbf{Developing all-atom protein JEPAs.} We adapt three predictive approaches to the La-Proteina encoder. JEPA predicts masked residue representations with a target encoder updated by an exponential moving average (EMA), JEPA+SIGReg uses a shared encoder and SIGReg regularisation, and LeJEPA aligns protein or residue representations across views. They allow us to examine different ways to learn from all-atom structures without reconstructing coordinates as the training target (Section~\ref{sec:methodology}).
    \item \textbf{Designing protein views and objectives.} We investigate random, spatial and chemical masking, positional information, independent rotations, and protein- versus residue-level alignment. For LeJEPA, we also develop fixed-length global and local crops. We study how these design choices change the latent space and affect downstream performance (Sections~\ref{sec:methods-views},~\ref{sec:methods-cropped-lejepa} and~\ref{sec:results-masking}).
    \item \textbf{Evaluating and probing learned representations.} We compare the predictive models with two reconstruction-based autoencoders and randomly initialised controls on twelve StructTokenBench tasks spanning protein-level remote homology, residue-level functional sites and flexibility. We use t-SNE visualisations and linear probes to study latent-space organisation. We also reconstruct all-atom structures from a selected JEPA representation, compare them with the La-Proteina baseline, and visualise complete proteins and local decoder failures (Sections~\ref{sec:methods-evaluation},~\ref{sec:results-transfer} and~\ref{sec:results-reconstruction}).
    \item \textbf{Interpreting model behaviour.} We examine how absolute residue position and chain length can become shortcuts, and introduce relative-position prediction and fixed-length crops to address them. We compare intermediate encoder blocks with the final 8D latent to locate where downstream-useful information is retained (Sections~\ref{sec:results-organisation} and~\ref{sec:results-layerwise}).
\end{enumerate}

\paragraph{Thesis outline.}
Section~\ref{sec:background} introduces protein structure, representation learning, the La-Proteina encoder and the JEPA objectives used in this thesis. Section~\ref{sec:related-work} reviews related protein foundation models, structural representation methods and protein JEPAs, and identifies the gap addressed by this work. Section~\ref{sec:methodology} presents the data, model adaptations, protein views and evaluation protocols. Section~\ref{sec:results} reports the experimental results and critically analyses downstream transfer, masking choices, reconstruction, latent-space organisation, representation collapse and intermediate encoder blocks. Finally, Section~\ref{sec:conclusion} summarises the contributions, discusses their limitations and proposes directions for future work.
